% 图 4.5 MNN 概述示意图 — 基于 CICC0903540 技术文档 4.3
\documentclass[tikz,border=8pt]{standalone}
% 顶会/顶刊风格：低饱和度马卡龙 + 高级感
% 适用于 NeurIPS, CVPR, ICLR, IEEE TPAMI
\usepackage{amssymb}
\usepackage{fontspec}
\usepackage{xeCJK}
\setCJKmainfont{SimHei}
\usepackage{tikz}
\usepackage{pgfplots}
\usetikzlibrary{
  arrows.meta, positioning, shapes.geometric, calc,
  backgrounds, fit, decorations.pathreplacing, patterns
}

% 低饱和度马卡龙配色
\definecolor{mBlue}{HTML}{AEC6CF}
\definecolor{mPink}{HTML}{FFD1DC}
\definecolor{mGreen}{HTML}{B1D8B7}
\definecolor{mPurple}{HTML}{B39EB5}
\definecolor{mYellow}{HTML}{FDFD96}
\definecolor{mGray}{HTML}{E5E4E2}
\definecolor{mDark}{HTML}{4A4A4A}

% 全局 TikZ 样式
\tikzset{
  >=stealth,
  every node/.style={align=center, font=\sffamily},
  block/.style={
    rounded corners=3pt, draw=mDark, align=center,
    minimum height=0.8cm, inner sep=5pt},
  arr/.style={-{Stealth[length=4pt]}, semithick, draw=mDark},
  % 信息密度增强：张量维度标注
  ts/.style={font=\tiny\sffamily, color=mDark!80},
  % 数学操作符节点（⊕加 ⊗乘 ‖拼接）
  op/.style={circle, draw=mDark, minimum size=0.4cm, inner sep=0,
    font=\scriptsize\sffamily, fill=mGray!60},
  % 图例
  legendbox/.style={rectangle, draw=mDark!50, rounded corners=2pt,
    fill=mGray!15, inner sep=4pt, font=\tiny\sffamily},
}

\begin{document}
\begin{tikzpicture}[
  x=1.0cm, y=1.0cm,
  node distance=0.5cm and 0.7cm,
  box/.style={
    rectangle, draw=mDark, fill=mBlue!40,
    minimum height=0.8cm, font=\small\sffamily, align=center,
    rounded corners=4pt, semithick, inner sep=5pt},
  subbox/.style={
    rectangle, draw=mDark!80, fill=mGray!50,
    minimum height=0.65cm, font=\scriptsize\sffamily, align=center,
    rounded corners=3pt},
  arr/.style={-{Stealth[length=4pt]}, semithick, draw=mDark},
  lbl/.style={font=\footnotesize\sffamily\bfseries, color=mDark},
]
  \node[box, minimum width=11cm] (input)
    {输入模型（TensorFlow / TFLite / Caffe / ONNX）};

  \node[lbl, below=0.8cm of input.west, anchor=west] (cvt) {Converter};
  \node[subbox, minimum width=4.8cm, below=0.25cm of cvt.south west, anchor=north west] (fe)
    {Frontend（格式解析）};
  \node[subbox, minimum width=4.8cm, right=0.5cm of fe] (go)
    {Graph Optimizer（算子融合）};

  \node[lbl, below=0.9cm of fe.south west, anchor=west] (int) {Interpreter};
  \node[subbox, minimum width=4.8cm, below=0.25cm of int.south west, anchor=north west] (eng)
    {Engine（调度/内存管理）};
  \node[subbox, minimum width=4.8cm, right=0.5cm of eng] (be)
    {Backend（CPU/GPU/NPU）};

  \node[subbox, minimum width=3.2cm, below=0.6cm of eng] (geo)
    {几何计算\\{\scriptsize\color{mDark!80} 光栅算子分解}};
  \node[subbox, minimum width=3.2cm, right=0.5cm of geo] (semi)
    {半自动搜索\\{\scriptsize\color{mDark!80} 最优后端选择}};
  \node[subbox, minimum width=2.6cm, right=0.5cm of semi] (neon)
    {NEON/汇编\\{\scriptsize\color{mDark!80} 底层加速}};

  \draw[arr] (input.south) -- ++(0,-0.3) -| (fe.north);
  \draw[arr] (fe) -- (go);
  \draw[arr] (go.south) -- ++(0,-0.25) -| (eng.north);
  \draw[arr] (eng) -- (be);
  \draw[arr] (eng.south) -- ++(0,-0.15) -| (geo.north);
  \draw[arr] (geo) -- (semi);
  \draw[arr] (semi) -- (neon);

  \begin{scope}[on background layer]
    \node[fit=(cvt)(fe)(go), fill=mBlue!25, draw=mDark!50,
      rounded corners=5pt, inner xsep=8pt, inner ysep=6pt, semithick] {};
    \node[fit=(int)(eng)(be)(geo)(semi)(neon), fill=mPurple!25, draw=mDark!50,
      rounded corners=5pt, inner xsep=8pt, inner ysep=6pt, semithick] {};
  \end{scope}

  % 图例（预留边距）
  \node[legendbox, anchor=north east] at ([xshift=-6pt, yshift=-6pt]current bounding box.south east)
    [align=left, text width=3.8cm] {
    \textbf{图例}\\[1pt]
    \colorbox{mBlue!40}{\rule{0.15cm}{0.1cm}\hspace{0.06cm}} Converter \quad
    \colorbox{mPurple!40}{\rule{0.15cm}{0.1cm}\hspace{0.06cm}} Interpreter\\
    61原子+45转换+16复合 $\to$ 光栅+原子 $O(1055)$
  };
\end{tikzpicture}
\end{document}
